\documentclass[sigconf]{acmart}

%\usepackage{booktabs} % For formal tables
%\usepackage[utf8]{inputenc}

% Copyright
\setcopyright{rightsretained}

\title{Armchair Auditing of Insolvency Processes}
\author{Tom Akkermans}
\affiliation{\institution{University of Amsterdam}}
\email{tom.akkermans@student.uva.nl}


\begin{document}
\begin{titlepage}
\begin{center}

\textsc{\Large Data Driven Supervision in Insolvency Procedures }

\bigskip

\textsc{\large
submitted in partial fulfillment for the degree of master of science\\
%
\bigskip
Tom Akkermans\\
%
11323671\\
%
\bigskip
master information studies\\
%
data science \\
%
faculty of science\\
%
university of amsterdam\\
%
\bigskip
%Your Date of defence in the format YYYY-MM-DD
2018-08-18
}

\end{center}
 

\vfill

% In case of an internal project, remove External Supervisor or if you had two internal supervisors, change the header into 
%  & First Supervisor & Second Supervisor  \\
\begin{center}
\begin{tabular}{|l||ll|}
\hline
 & \textbf{First Supervisor} & \textbf{Second Supervisor}  \\   
 \hline
\textbf{Title, Name} & Dr Maarten Marx&  \\
\textbf{Affiliation} &UvA, FNWI, IvI & \\ 
\textbf{Email} & maartenmarx@uva.nl& . \\
\hline
\end{tabular}
\end{center}

\bigskip

% logos
\begin{center}
\mbox{\includegraphics[width=.2\paperwidth]{logo-uva.png} 
\includegraphics[width=.2\paperwidth]{ads.png}
}
\end{center}
\end{titlepage}

%
%\newpage
%
%\end{document}


\begin{abstract}
When a company is declared bankrupt at the court, a court committee will appoint a curator to settle the bankruptcy. The curator's task is to redeem the creditors claims as much as possible by monetizing the company's estate. A court's judge supervises the curator to act in the best interest of the creditors.

The curator appointment process, the judge supervision capability and the curator-creditor interest alignment have been topic of research[[1-2]], raised concerns in the media \cite{dennis_meneer_2018:1} and have lead to court cases for not being transparent and/or not being adequate. Involved parties such as judges and curators as well as the creditors and insolvent itself would benefit from more transparency and information on these processes.

The Dutch government has initiated a transparency program for the legal system. It provides on-line search form for a single insolvency case and provides open data sources with technical APIs supplying court publications and curator reports. But the search form is limited and not everyone with an interest can handle the raw data of the APIs. Instead of open data, the stakeholder needs open analysis to enable armchair auditing of insolvency processes to answer their questions.

In this thesis an information system prototype will be developed to enable the auditing of insolvency processes for non tech savvy users. The system takes in structured and unstructured data of insolvency procedures using open and publicly available data sources[[22-28]]. From this data the system builds and annotates a complete linked entity structure and unstructured data is made searchable on multiple levels. The stakeholder can query the linked data with a non-technical interface in order to explore insolvency processes and answer specific questions.
\end{abstract}

%
% The code below should be generated by the tool at
% http://dl.acm.org/ccs.cfm
% Please copy and paste the code instead of the example below.
%
\begin{CCSXML}
<ccs2012>
	<concept>
		<concept_id>10002951.10003317.10003347.10003349</concept_id>
		<concept_desc>Information systems~Document filtering</concept_desc>
		<concept_significance>300</concept_significance>
	</concept>
</ccs2012>
\end{CCSXML}
\ccsdesc[300]{Information systems~Document filtering}

\keywords{keyword1, keyword2, keyword3}

\maketitle 


\section{Introduction}
When a company is declared bankrupt by the court, a court committee appoints an administrator to settle the bankruptcy. The administrator's task is to liquidate the company's estate and use the proceeds to settle the creditors claims. A supervisory judge ensures that the administrator is acting in the best interest of the creditors. 

Involved parties have been demanding more transparency into insolvency processes. The supervisory function of the judge, the conflict of interest between the administrator and creditors and the appointment process of the administrator are processes which have been the subject of research\cite{boluk_2011}, about which media articles have appeared \cite{dennis_meneer_2018:1, dennis_meneer_2017:1, jan-hein_strop_2015:1} and which have led to legal proceedings. \todo{citeer}  

\begin{comment}
Other benefits: Supervisory judges with significant work load could benefit from data driven supervision. Information access to the general public and journalists to this processes  would provide additional checks and balances.
\end{comment}

In 2005 the Dutch government started the digital register of insolvency data\cite{rechtspraak:cir}. An on-line search form \cite{rechtspraak:cir-zoeken} is provided to retrieve a single insolvency case. Web services are provided to retrieve court publications in XML format and administrator reports in PDF format. 

However, the information from a single insolvency case is limited as it does not provide aggregated and linked information. The administrator reports are unstructured and not searchable as a collection. Furthermore most involved parties lack the technical skills to use the web services. Instead of open data, there is a need for \textbf{open analysis} to enable 'armchair audits'\cite{o_leary_2015} of insolvency processes.

% main research question
In this thesis we investigate whether [RQ] \textbf{it is possible to build a complete and correct structured information system (IS) based on provided and enriched open and public data that is able to provide the requested transparency to the parties involved}.

For this we need to answer the following questions:
\begin{itemize}
	\item [RQ1] Can the IS construct a complete, cleaned and fully linked entity network of insolvency cases, administrators, judges and courts.
	\item [RQ2] Can the IS correctly and completely label insolvency case data with state data [start/end date] in order to mine the insolvency process.
	\item [RQ3] Can the IS correctly extract specific fields [paulianeus handelen] of interest from unstructured documents to classify insolvency cases.
	\item [RQ4] Can involved stakeholders use the IS to provide the requested transparency on insolvency processes.
\end{itemize}

We describe the steps in building such a system that takes in large amounts of open and publicly available data in structured and unstructured form, extracts and enriches useful facts and makes it consumable for analysis to provide insights into the insolvency processes via a web GUI.

\section{Related Work}
\todo[inline]{add related work on entity linking, deduplication etc}

\subsection{Research Questions}
The main research question is:
\textbf{Is is possible to build a useful, complete and correct structured information system based on unstructured CIR data?}

The system is build on three components
\begin{itemize}
\item extraction of insolvency process flow information.
\item construction of a fully linked, clean entity structure of insolvents, administrators, judges, courts as well as administrator reports and court publications.
\item extraction the text and parameters in administrator reports and indexes sections and parameters by imposing structure on the content.
\end{itemize} 

\begin{itemize}
\item \textbf{useful}: the system is useful when it can answer questions of the parties involved or can answer these questions much more efficient than with current methods. We define so-called Persona to represent archetypical users of the information system and define specific questions they have. These questions are distilled from the insolvency law, news articles, research papers, court cases [refs] and interviews. [How do we measure, which section]
\item \textbf{complete} the system's entity structure is complete when all main entities and their inter relations are found. Completeness is also defined for specific parameters needed to answer the user's questions. [How do we measure, which section]
\item \textbf{accuracy} The system's identity structure is accurate when entities refer to their actual real life counterparts. The [How do we measure, which section]


\end{itemize}


\section{Methodology}
\input{chapters/methodology-data-sources}

\paragraph{Database normalization}
This leads to unwanted duplication as names can be written in many different forms and typos can be introduced, e.g. one judge's name appears as:

\begin{itemize}
\item "mr. W.J. Geurts - de Veld"
\item "mr. W.J. Geurts-deVeld"
\item "mr. W.J. Geurts-de Veld"
\item "mr. W.J.Geurts-de Veld"
\item "mr.W.J. Geurts-de Veld"
\item "mr. W.J. Geurts-de Veld (Rotterdam)"
\item "mr W.J.Geurts-de Veld"
\item "W.J.Geurts-de Veld"
\end{itemize}

We need to define so-called master data for judges and administrators containing the real world entities. The two data sources in the sections \ref{NOvA Tableau} and \ref{Nevenfuncties Rechters} are chosen for this purpose. CIR entity names are first normalized for de-duplication and are subsequently linked to the master data records on their normalized name. 

\paragraph{PDF report web service}




The number of current defaults is about the lowest of the century [ref] [graph of declining defaults]






\subsection{Implementation challenges}
\subsubsection{Entity De-duplication, Scraping and Linking}
\paragraph{Deduplication: normalizing names}
CIR provides administrator names in four parts: title, initials, middle part and family name. 

\textbf{The initials} were normalized into single characters using periods and no spaces. In Dutch, initials for names like \textit{Theodoor} and \textit{Christiaan} are written with two or three characters as \textit{Th.} or \textit{Chr.}. These names are derived from the Greek where the initial \textit{Th(eta)} for example is written as one character. Initials for double names like \textit{Albert-Jan} may be written as \textit{A-J.} or \textit{A.-J.}. 

\textbf{The middle part}, e.g. 'van der' or 'ter' was sometimes supplied in the middle part field and sometimes in the family name (and sometimes in both). To normalize the middle part it was merged into the family name making sure it did not appear twice.

\textbf{The family name} was stripped from academic and noble (not uncommon in this profession) titles. Misplaced initials and CIR specific comments were removed. Maiden names are connected to the surname with a hyphen and no spaces.

Furthermore, all parts are made lower case, surrounding whitespace characters are stripped, accents are removed and multiple spaces are replaces with a single space. 

De-duplication by normalizing the administrator names reduced the number of distinct administrator names from 2329 to 1559, a 33\% reduction.
\\\\
Judge names are provided in a single field. This sometimes leads to initials or a middle name part that is stuck together to the family name and was separated using regexm e.g. in 'C. vanSteenderen' and 'mr. W.J.Geurts-de Veld'. CIR often adds the court between parenthesis which was also removed in the normalization.

De-duplication by normalizing the judge names reduced the number of distinct judge names from 565 to 277, a 51\% reduction.

\paragraph{Administrator name scraping}
NOvA unfortunately does not provide complete lists of registered lawyers to which we could link the administrators. Instead we link the administrator by using the site's search form. An underlying REST API returning JSON data was discovered which was used instead of the form. The API (or form) can not handle initials or stop words as 'de' and 'van' and only the family name without stop words was used. The best candidate in the search results, which are sorted on the relevance provided by NOvA, is found by (partially) matching the initials until a match is found. The API enables the use of filters and we filter from narrow to broad until a candidate is found, filtering on: 
\begin{enumerate}
\item{legal specialism: administrator}
\item{legal branch: insolvency law}
\item{no filter}
\end{enumerate}

NOvA only supplies the actual lawyer register but CIR also lists administrators previously working on a case therefore a substantial amount of administrators cannot be linked. Another website, www.advocatenzoeken.nl that sources data from NOvA but keeps historical data was scraped as well to complement NOvA.

\paragraph{Administrator name linking}
From the normalized name list we linked 73\% to NOVa registered lawyers and another 18\% administrators from the www.advocatenzoeken.nl site. The remaining 10\% unfound names are linked to a special 'unknown' lawyer.

\begin{table}[h]
\caption{Administator entity linking results}
\centering
\begin{tabular}{l r r r}
\hline\hline
Source & no. linked & correct & incorrect\\
\hline
NOvA & 1134 & 1133 & 1\\
AdvocatenZoeken & 280 & 279 & 1 \\
Not found & 149 &&\\
\hline
Total & 1563 & 1412 &\\
\hline
\end{tabular}
\label{table:administrator_linking}
\end{table}

Correctness is assumed when the names match. When an obvious typo is found in the family name. When the initials either match, or are contained in the master record or are permutated or are missing. 

[discuss not found results: common name, place in name]

\paragraph{Judge name scraping}
The CIR register for ancillary functions for judges was scraped for judges from courts and higher courts. The total set contains 3691 judges and should be a superset of active case judges. The website was driven by the Angular javascript framework and could not easily be scraped with simple HTTP requests. Selenium was used to mimic user browsing behaviour to invoke the javascript. 

\paragraph{Judge entity linking}
A judge is linked by to a judge master record with the most similar normalized name where similarity was calculated using the Levensthein or edit distance. 88\% of the normalized names were matched. 4 of the incorrectly 12\% matched names could be set manually and the remaining 8\% was linked to a special 'unknown' judge.


\begin{table}[h]
\caption{Judge entity linking results}
\centering
\begin{tabular}{l r r r}
\hline\hline
Source & no. linked & correct & incorrect\\
\hline
Ancillary positions register & 277 & 246 & 32\\
+Manual lookup of the 32 &   & 10 & -10\\
\hline
Total & 277 & 256 & 22\\
\hline
\end{tabular}
\label{table:judge_linking}
\end{table}

Correctness is defined similar to the administrator normalized name correctness.

[discuss not found (initially) results: maiden name, not a judge anymore]

\subsubsection{Process Mining}

\subsubsection{PDF content extraction}
The PDF reports can be split into two types depending on how the PDF was created:
\begin{enumerate}
\item Scanned from paper using a scanner or copier.
\item Converted by software from another format.
\end{enumerate}

\paragraph{Scanned PDFs}
Scanned PDFs contain images only. To convert the images to text we used the Ghostscript library for image extraction and the Tesseract OCR engine for the character recognition. Tesseract supports the Dutch language which is paramount to our application. It is very important to supply Tesseract with good quality images. We used the Ghostscript settings of a \emph{tiffgray} output device with a \emph{300 DPI} resolution. 

\paragraph{Converted PDFs}
Converted PDF content can be extracted with a number of packages such as PDFMiner, PyPDF2, pyPoppler and pdftotext. We chose the latter after comparison because it maintains the layout which is important for section and parameter extraction and being build in C++ it is very fast.

\paragraph{Third type}
A third type appeared: PDFs that were scanned and subsequently OCR-ed by a copier. They contain both text and images. The OCR quality is often poor. We retried re-OCR-ing the images with Tesseract which solved some errors but introduced others. Meta data about a.o. the scanner type was obtained which could improve post-process text extraction in future work.


\subsubsection{PDF report section and parameter extraction}
[todo]




\input{entity_resolution}

\subsection{Methods}
\subsubsection{RQ1}
\section{Evaluation}
dedup lawyers: Phonetic similarity
https://www.quora.com/What-is-a-good-algorithm-service-for-fuzzy-matching-of-companies-names-for-de-duplication


\section{Conclusions}
\subsection{Acknowledgements}



\bibliographystyle{ACM-Reference-Format}
\bibliography{bibliography}

\end{document}
